\documentclass[12pt]{article}
\usepackage[utf8]{inputenc}
\usepackage[T1]{fontenc}
\usepackage{hyperref}
\title{Phowa, Hlepho Language Resources for K-12 Curriculum Development}
\author{Compiled by Claude Sonnet 5.5 (claude-sonnet-5-5) for Andr\'as Kornai}
\date{}
\begin{document}
\maketitle

\section{Introduction}
Hlepho Phowa (ISO 639-3 yhl; Glottocode hlep1236) is a Loloish (Ngwi) Sino-Tibetan language of Yunnan, China. Ethnologue's 2011 figure is about 36,000 speakers. Wikitongues rates it 6b (Threatened). It is not known to be taught in schools. The search found only reference and overview pages, and no educational materials.

\section{Primary Educational Resources}
No primers, textbooks or school programmes were found.

\section{Linguistic and Reference Materials}
\begin{itemize}
\item Phowa language (Wikipedia). It covers classification and varieties (Hlepho, Ani, Labo). \url{https://en.wikipedia.org/wiki/Phowa_language}
\item Pelkey, ``Puzzling over Phula'' (Academia.edu). It describes Phowa as a newly documented Phula language. \url{https://www.academia.edu/2435607/Puzzling_over_Phula}
\item Language name and location: Hlepho Phowa, Yunnan (MPI EVA numerals page). \url{https://lingweb.eva.mpg.de/channumerals/Hlepho-Phowa.htm}
\end{itemize}

\section{Religious and Community Resources}
\begin{itemize}
\item Joshua Project, Phowa, Hlepho in China. It lists the language as needing a Bible translation. \url{https://joshuaproject.net/people_groups/18657}
\end{itemize}

\section{Digital and Audio Resources}
\begin{itemize}
\item Wikitongues, Hlepho Phowa. This is a language profile page. Whether it holds audio or video was not confirmed. \url{https://wikitongues.org/languages/yhl/}
\end{itemize}

\section{Curriculum Development Context}
Most adults speak the language, but not all young people do, and it is not known to be taught in schools. Pelkey's work is the main documentation found. A curriculum would need to build primary materials, such as an orthography and a primer, from scratch. Further searching, for example in Chinese-language sources, is needed.

\end{document}